%------------------------------------------------------------------------------%
\begin{question}
  Considerando a função
  \(
    f(x, y) = x^4 + y^4 + 4xy
  \)

  a) Calcule o gradiente de $f$

  b) Calcule a hessiana de $f$

  c) Encontre todos os pontos críticos de $f$

  d) Classifique cada ponto crítico de $f$
\end{question}

%------------------------------------------------------------------------------%
\begin{resolution}{a}

  Precisamos das derivadas parciais de $f$
  \[
    \D{f}{x}
    = \D{}{x}\left[x^4 + y^4 + 4xy\right]
    = 4x^3 + 4y
  \]
  \[
    \D{f}{y}
    = \D{}{y}\left[x^4 + y^4 + 4xy\right]
    = 4y^3 + 4x
  \]
  então
  \[
    \nabla f = \vetor{4x^3 + 4y}{4y^3 + 4x}
  \]
\end{resolution}

%------------------------------------------------------------------------------%
\begin{resolution}{b}

  Precisamos das derivadas parciais de segunda ordem de $f$
  \[
    \DS{f}{x}
    = \D{}{x}\left[4x^3 + 4y\right]
    = 12x^2
  \]
  \[
    \D{f}{y}
    = \D{}{y}\left[4y^3 + 4x\right]
    = 12y^2
  \]
  \[
    \DM{f}{x}{y}
    = \D{}{x}\D{f}{y}
    = \D{}{x}\left[4y^3 + 4x\right]
    = 4
  \]
  então
  \[
    H = \left(\begin{array}{cc}
      12x^2 & 4 \\
      4 & 12y^2
    \end{array}\right)
  \]
\end{resolution}

%------------------------------------------------------------------------------%
\begin{resolution}{c}
  A função é um polinômio, então possui derivadas em todos os pontos do plano.
  Assim os pontos críticos são apenas os pontos onde as derivadas parciais são zero,
  \(\nabla f = 0\)
  \[
    4x^3 + 4y = 0
    \qquad\text{e}\qquad
    4y^3 + 4x = 0
  \]
  ou, simplificando,
  \[
    x^3 + y = 0
    \qquad\text{e}\qquad
    y^3 + x = 0
  \]
\end{resolution}

%------------------------------------------------------------------------------%
\begin{resolution}{c}
  Isolando $y$ na primeira equação, \(y = -x^3\), e substituindo na segunda, temos
  \begin{align*}
    y^3 + x                 & = 0 \\
    \left(-x^3\right)^3 + x & = 0 \\
    -x^9 + x                & = 0 \\
    x^9 - x                 & = 0 \\
    x(x^8 - 1)              & = 0
  \end{align*}
  As soluções dessa equação são $x=0$, $x=1$ ou $x-1$.
  Se $x=0$ temos $y=0$,
  se $x=1$ temos $y=-1$ e
  se $x=-1$ temos $y=1$.
  Portanto, os pontos críticos são
  \[
    (x_1, y_1) = (0, 0),
    \qquad\qquad
    (x_2, y_2) = (1, -1)
    \qquad\text{e}\qquad
    (x_3, y_3) = (-1, 1)
  \]
\end{resolution}

%------------------------------------------------------------------------------%
\begin{resolution}{d}
  Para classificar os pontos críticos precisamos avaliar o discriminante nos
  pontos críticos.
\end{resolution}

%------------------------------------------------------------------------------%
\begin{resolution}{d}
  Considerando o ponto \((x_1, y_1) = (0, 0)\)
  \[
  f_{xx}(0, 0) = 0
  \]
  \[
  f_{yy}(0, 0) = 0
  \]
  \[
  f_{xy}(0, 0) = 4
  \]
  \[
  D_1
  = f_{xx}(0, 0)f_{yy}(0, 0) - f^2_{xy}(0, 0)
  = 0 \times 0 - 4^2
  = -16 < 0
  \]
  Portanto, o ponto $(0, 0)$ é um ponto de sela.
\end{resolution}

%------------------------------------------------------------------------------%
\begin{resolution}{d}
  Considerando o ponto \((x_2, y_2) = (1, -1)\)
  \[
    f_{xx}(1, -1) = 12
  \]
  \[
    f_{yy}(1, -1) = 12
  \]
  \[
    f_{xy}(1, -1) = 4
  \]
  \[
    D_2
    = f_{xx}(1, -1)f_{yy}(1, -1) - f^2_{xy}(1, -1)
    = 12 \times 12 - 4^2
    = 144 -16
    = 128 > 0
  \]
  Portanto, o ponto $(1, -1)$ é um máximo ou mínimo local.
  Como \(f_{xx}(1, -1) = 12 > 0\)
  o ponto é um ponto de mínimo local.
\end{resolution}

%------------------------------------------------------------------------------%
\begin{resolution}{d}
  Considerando o ponto \((x_3, y_3) = (-1, 1)\)
  \[
    f_{xx}(-1, 1) = 12
  \]
  \[
    f_{yy}(-1, 1) = 12
  \]
  \[
    f_{xy}(-1, 1) = 4
  \]
  \[
    D_3
    = f_{xx}(-1, 1)f_{yy}(-1, 1) - f^2_{xy}(-1, 1)
    = 12 \times 12 - 4^2
    = 144 -16
    = 128 > 0
  \]
  Portanto, o ponto $(-1, 1)$ é um máximo ou mínimo local.
  Como \(f_{xx}(-1, 1) = 12 > 0\)
  o ponto é um ponto de mínimo local.
\end{resolution}

%------------------------------------------------------------------------------%
